% Pursuing Stacks (cotes 134-2 to 134-8): archive cote 134-4, site page 26 (PDF page 27), typescript page 114.
% typescript. Draft compared: ps3.tex:1281-1345.
% AI-transcribed by Claude Opus 5.5 (claude-opus-5-5, adaptive thinking, effort high), Message Batches API, batch ps2 (msgbatch_01NVwr1Yuo8qJf5rykeGBXPK), 2026-11-12; image review recorded in records/CHECK_DECISIONS_134_3_4.tsv.
\scanpage{134-4/26}
\mspage{114}

as a \uline{basic assumption} after all.  This assumption makes sense of course
independently of any saturation condition.  If $W_h$ is not \xout{\unsure{satura}} strongly
saturated,  then denoting by $\overset{\hand{\rule{0.6em}{0.4pt}}}{W}_h$ the subset of $F\over{l}{l}(M)$ of all arrows made
invertible by the canonical functor,  this will now be a strongly satura-%
ted set of arrows (in fact the smallest one containing $\over{W}{W}_h$\deleted{)},  and giving
rise to the same localise\over[typed]{\illegible{1}}{d} category,  and hence satisfying the basic \xout{\unsure{assum}}
assuption too\hand{)} -  thus $\overset{\hand{\rule{0.6em}{0.4pt}}}{W}_h$ is indeed a homotopism structure on M.

When M admits a final object,  this will equally be a final object of
$W_h^{-1}M$ . We may now define in terms of $W_h$  the notion of contractible
\xout{elements} objects in M , \over{\illegible{1}}{forming} a subset
\[ C_h \hand{\subset} Ob(M) \quad , \]
as those objects X in M such that  the projection $p_X : X \hand{\ensuremath{\to}} e$ is in $W_h$ ,
or , equivalently,  such that X is \xout{isomorphic} a final object in the loca-%
lized category $W_h^{-1}M$ .  This accounts for the th\over{\unsure{ri}}{ir}d plain arrow of the
\uline{organigram.}\diagram{\begin{tikzpicture}[baseline=-0.4ex,x=1em,y=1ex]\draw (0,0)--(0.4,0)--(0.4,1.6)--(1.8,1.6);\end{tikzpicture}} We'll now \xout{\illegible{1}}\inserted{dwell} a little more on the first dotted arrow,
namely the description of a homotopy relation
\[ R_h \hand{\subset} F\over{l}{l}(M)\mathrm{x}F\over{l}{l}(M) \]
in terms of $W_h$ :  \over{T}{t}\over[typed]{H}{h}e natural choice here is to define  f,g\hand{$\in$} \xout{xxxxx}\inserted[typed]{$F\over{l}{l}(M)$} to
be \uline{homotopic} (\uline{or $W_h$}-homotopic, if ambiguity may arise)  iff their images
in the category $W_h^{-1}M$  are equa\over{\illegible{1}}{l}.  This \xout{homotopy notion}\inserted[typed]{relation} between maps
in M clearly satisfies the basic assumption on source and target,  as well
as the saturation condition  -  it is  therefore a``homotopy \over{\illegible{1}}{rela}tion'' in M,
\xout{and} namely the one associated to $\over{W}{W}_h$\hand{$^{-1}M$,} viewed as a quotient category of M.
It is clear that we recover $W_h$ from $R_h$ ,  consequently,  by the process
described in A),  \xout{of taking}

To make the relationship between the notions 1\hand{)} and 2\hand{)} still clearer,
let\hand{'}s denote respectively by
\[ Ho\over[typed]{M}{m}_1(M),\ Hom_2(M) \]
the set of all homotopy \xout{\unsure{notions resp.}} relations,  resp. of all homotopism
notions,  in M.  We got maps\hand{*}
% Hand-drawn arrows are marked blue; typed objects remain plain.
\diagram{\begin{tikzcd} Hom_1(M) \arrow[r, shift left, color=blue, "\hand{r_{21}}"] \& Hom_2(M) \arrow[l, shift left, color=blue, "\hand{r_{12}}"] \end{tikzcd} \quad ,}
\margin{* 3.4. There \uline{is no} \uline{map} $r_{21}$, only $r_{12}$, see correction in \S\ 51!}
and the relevant fact here is that
\[ r_{12}: Hom_2 \hand{\ensuremath{\longrightarrow}} Hom_1 \]
\uline{is injective},  \uline{and admits} $r_{21}$ \uline{as a left inverse}.  Thus we may view
$Hom_2$ as a subset of $Hom_1$ ,  i.e. the structure of a ``homotopism notion'' on
M as a particular case of the structure of a ``homotopy relation'' on M.
Namely,  a structure of the latter type can be described in terms of a
notion of homotopism in M,  iff the canonical functor $M \hand{\ensuremath{\to}} M_h$ i\over{s}{t} gives rise
\xout{to}

\margin{26}
